We now extend the result to more general domains.

\begin{Theorem}
\label{thm:smooth-domain-sobolev-extension}
Assume that $\Omega$ is of class $C^K$ and $\partial\Omega$ is bounded. Then, for any integer $k$ with $0\leq k\leq K$, there exists a bounded operator\[E:H^k(\Omega)\to H^k(\mathbb{R}^m)\]such that $Eu|_{\Omega}=u$. Moreover, $E$ can be chosen independent of $k$.
\end{Theorem}

Note that we only require $\partial\Omega$ to be bounded, not necessarily $\Omega$. For example, the theorem includes the case of exterior domains.

\begingroup
\postdisplaypenalty=10000
\renewcommand{\endProof}{\par\nobreak\hfill$\blacksquare$\par\vspace{0.1cm}}
\begin{Proof}
We cover $\partial\Omega$ with small neighborhoods $U_j$ such that within each $U_j$ there is a $k$-diffeomorphism which maps $\partial\Omega\cap U_j$ to a subset of $\{x_m=0\}$. Let $\{\phi_j\}$ be a partition of unity subordinate to the covering $\{U_j\}$ of $\partial\Omega$. Let $u\in H^k(\Omega)$. Then, for each $j$, $\phi_ju$ can be extended to all of $\mathbb{R}^m$ by using the transformation theorem and Theorem~\ref{thm:halfspace-sobolev-bounded-extension}. More precisely, let $\bm{\Phi}_j$ be the diffeomorphism on $U_j$ which maps $\partial\Omega\cap U_j$ to a part of the plane $x_m=0$. Then $(\bm{\Phi}_j^{-1})^*(u\phi_j)$ (extended by zero outside $\bm{\Phi}_j(U_j)$) is in $H^k(\mathbb{R}_+^m)$ and can be extended to a function in $H^k(\mathbb{R}^m)$; let us call the extended function $v_j$. We multiply $v_j$ by a function $\psi_j\in\mathcal{D}(\mathbb{R}^m)$ which has support in $\bm{\Phi}_j(U_j)$, but equals one in $\bm{\Phi}_j(\operatorname{supp}\phi_j)$. Then $\bm{\Phi}_j^*(v_j\psi_j)$ is in $H^k(\mathbb{R}^m)$ with support contained in $U_j$. Hence each $\phi_ju$ can be extended. Finally $u-\sum_j\phi_ju$ has compact support in $\Omega$, hence it can be extended by zero (see Problem~\ref{ex:compact-support-zero-boundary-sobolev}).
\end{Proof}
\endgroup

\begin{Remark}
\label{rem:lipschitz-domain-sobolev-extension}
The same result (with a much harder proof) holds if $\partial\Omega$ is Lipschitz rather than of class $C^K$.
\end{Remark}
